\documentclass[a4paper,titlepage,12pt]{article}
\usepackage[utf8]{inputenc}
\usepackage[T1]{fontenc}
\usepackage[margin=1in]{geometry}
\usepackage{parskip}
\usepackage{graphicx}
\usepackage{hyperref}
\usepackage{listings}
\usepackage{multirow}
\usepackage[usenames,dvipsnames]{color}
\hypersetup{
	colorlinks,
	pdfauthor=Delan Azabani,
	pdftitle=Design and Analysis of Algorithms 300: Huffman assignment
}
\lstset{basicstyle=\ttfamily, basewidth=0.5em}

\title{Design and Analysis of Algorithms 300\\Huffman assignment}
\date{April 19, 2014}
\author{Delan Azabani}

\begin{document}

\maketitle

\section{Design considerations regarding data types}

One of my goals when completing this assignment was to ensure that the core
Huffman and Base64 algorithms only ever interact with sequences of octets
(\texttt{byte[]}). Beyond making them agnostic of text encodings, this would
allow them to be used with other data which may not necessarily be textual.

Of course, the implementation still needed to conform to the specification, and
behave as expected when tested with copied textual input. This led to the need
to make careful decisions as to where and how the interactions between strings
and octets are made.

The .NET Framework defines \texttt{System.Char} as a `UTF-16 code unit', which
is obviously not an octet. Since the release of Unicode 2.0 in 1996 however,
these 16-bit values no longer necessarily represent entire Unicode code points
either. For instance:

\begin{tabular}{|lllll|}
	\hline
	Glyph &
		Unicode code point &
		UTF-32 &
		UTF-16 \texttt{char} &
		UTF-8 \\\hline
	\raisebox{-3pt}{
		\includegraphics[height=12pt]{U+000041.png}
	} &
		\texttt{U+000041} &
		\texttt{00000041} &
		\texttt{0041} &
		\texttt{41} \\\hline
	\raisebox{-3pt}{
		\includegraphics[height=12pt]{U+0000A9.png}
	} &
		\texttt{U+0000A9} &
		\texttt{000000A9} &
		\texttt{00A9} &
		\texttt{C2 A9} \\\hline
	\raisebox{-3pt}{
		\includegraphics[height=12pt]{U+002202.png}
	} &
		\texttt{U+002202} &
		\texttt{00002202} &
		\texttt{2202} &
		\texttt{E2 88 82} \\\hline
	\raisebox{-3pt}{
		\includegraphics[height=12pt]{U+01F530.png}
	} &
		\texttt{U+01F530} &
		\texttt{0001F530} &
		\texttt{D83D DD30} &
		\texttt{F0 9F 94 B0} \\\hline
\end{tabular}

The code units in the UTF-16 and UTF-32 encodings need to have a defined
\emph{endianness} when serialised into a sequence of octets. To avoid this, I
chose UTF-8 as the internal representation of any text where possible, only
using \texttt{string} and \texttt{char} when dealing with the UI. UTF-8 also
provides the unique benefit of backwards compatibility with ASCII.

This raises an interesting problem, even though I implemented § 2.2 of the
assignment, where frequency tables with arbitrary-length symbols are
\emph{accepted}. For example, I would accept a frequency table with a symbol
consisting of the octets \texttt{F0 9F 94 B0}.

How would I \emph{generate} a useful frequency table for text, if all I am
dealing with are sequences of octets? The approach used for binary data would
yield four separate entries for the character \includegraphics[height=12pt]
{U+01F530.png} --- \texttt{F0}, \texttt{9F}, \texttt{94} and \texttt{B0}, which
is clearly a poor result. My solution, albeit inelegant, was to write two
frequency table generators:

\begin{itemize}
	\item
		one which deals purely in octets, for binary data
	\item
		one which is aware that the octets are a UTF-8 sequence
\end{itemize}

This is still better than a naïve approach where all of the algorithms use
\texttt{char} values. Combining such a configuration with a character like
\includegraphics[height=12pt]{U+01F530.png} would lead to disastrous results.
While \texttt{D83D} and \texttt{DD30} are not valid characters, they would
receive separate entries in the frequency table, causing the table to fail to
correctly display in the UI.

\end{document}
